\begin{ex}
Rút gọn phân số $\frac{6}{8}$ về phân số tối giản.
\shortans{3/4}
\loigiai{
$\frac{6 : 2}{8 : 2} = \frac{3}{4}$.
}
\end{ex}

\begin{ex}
Rút gọn phân số $\frac{9}{12}$ về phân số tối giản.
\shortans{3/4}
\loigiai{
$\frac{9 : 3}{12 : 3} = \frac{3}{4}$.
}
\end{ex}

\begin{ex}
Rút gọn phân số $\frac{10}{15}$ về phân số tối giản.
\shortans{2/3}
\loigiai{
$\frac{10 : 5}{15 : 5} = \frac{2}{3}$.
}
\end{ex}

\begin{ex}
Rút gọn phân số $\frac{14}{21}$ về phân số tối giản.
\shortans{2/3}
\loigiai{
$\frac{14 : 7}{21 : 7} = \frac{2}{3}$.
}
\end{ex}

\begin{ex}
Rút gọn phân số $\frac{15}{20}$ về phân số tối giản.
\shortans{3/4}
\loigiai{
$\frac{15 : 5}{20 : 5} = \frac{3}{4}$.
}
\end{ex}

\begin{ex}
Rút gọn phân số $\frac{18}{27}$ về phân số tối giản.
\shortans{2/3}
\loigiai{
$\frac{18 : 9}{27 : 9} = \frac{2}{3}$.
}
\end{ex}

\begin{ex}
Rút gọn phân số $\frac{24}{32}$ về phân số tối giản.
\shortans{3/4}
\loigiai{
$\frac{24 : 8}{32 : 8} = \frac{3}{4}$.
}
\end{ex}

\begin{ex}
Rút gọn phân số $\frac{25}{100}$ về phân số tối giản.
\shortans{1/4}
\loigiai{
$\frac{25 : 25}{100 : 25} = \frac{1}{4}$.
}
\end{ex}

\begin{ex}
Rút gọn phân số $\frac{50}{100}$ về phân số tối giản.
\shortans{1/2}
\loigiai{
$\frac{50 : 50}{100 : 50} = \frac{1}{2}$.
}
\end{ex}

\begin{ex}
Rút gọn phân số $\frac{75}{100}$ về phân số tối giản.
\shortans{3/4}
\loigiai{
$\frac{75 : 25}{100 : 25} = \frac{3}{4}$.
}
\end{ex}

\begin{ex}
Rút gọn phân số $\frac{16}{24}$ về phân số tối giản.
\shortans{2/3}
\loigiai{
$\frac{16 : 8}{24 : 8} = \frac{2}{3}$.
}
\end{ex}

\begin{ex}
Rút gọn phân số $\frac{30}{45}$ về phân số tối giản.
\shortans{2/3}
\loigiai{
$\frac{30 : 15}{45 : 15} = \frac{2}{3}$.
}
\end{ex}

\begin{ex}
Rút gọn phân số $\frac{36}{48}$ về phân số tối giản.
\shortans{3/4}
\loigiai{
$\frac{36 : 12}{48 : 12} = \frac{3}{4}$.
}
\end{ex}

\begin{ex}
Rút gọn phân số $\frac{42}{56}$ về phân số tối giản.
\shortans{3/4}
\loigiai{
$\frac{42 : 14}{56 : 14} = \frac{3}{4}$.
}
\end{ex}

\begin{ex}
Phân số $\frac{5}{7}$ có phải là phân số tối giản không? (Phải ghi 1, không ghi 0).
\shortans{1}
\loigiai{
Phải vì tử và mẫu không cùng chia hết cho số nào lớn hơn 1.
}
\end{ex}

\begin{ex}
Phân số $\frac{6}{9}$ có phải là phân số tối giản không? (Phải ghi 1, không ghi 0).
\shortans{0}
\loigiai{
Không phải vì cùng chia hết cho 3.
}
\end{ex}

\begin{ex}
Trong các phân số $\frac{2}{3}, \frac{4}{6}, \frac{5}{10}, \frac{7}{9}$, có bao nhiêu phân số tối giản?
\shortans{2}
\loigiai{
$\frac{2}{3}$ và $\frac{7}{9}$ (có 2 phân số).
}
\end{ex}

\begin{ex}
Tử số của phân số tối giản của $\frac{12}{18}$ là bao nhiêu?
\shortans{2}
\loigiai{
$\frac{12}{18} = \frac{2}{3}$, tử số là 2.
}
\end{ex}

\begin{ex}
Mẫu số của phân số tối giản của $\frac{12}{18}$ là bao nhiêu?
\shortans{3}
\loigiai{
Mẫu số là 3.
}
\end{ex}

\begin{ex}
Rút gọn phân số $\frac{40}{60}$ được phân số tối giản là bao nhiêu?
\shortans{2/3}
\loigiai{
$\frac{40}{60} = \frac{2}{3}$.
}
\end{ex}

\begin{ex}
Rút gọn phân số $\frac{35}{49}$ được phân số tối giản là bao nhiêu?
\shortans{5/7}
\loigiai{
$\frac{35 : 7}{49 : 7} = \frac{5}{7}$.
}
\end{ex}

\begin{ex}
Rút gọn phân số $\frac{64}{80}$ về tối giản.
\shortans{4/5}
\loigiai{
$\frac{64 : 16}{80 : 16} = \frac{4}{5}$.
}
\end{ex}

\begin{ex}
Rút gọn phân số $\frac{81}{90}$ về tối giản.
\shortans{9/10}
\loigiai{
$\frac{81 : 9}{90 : 9} = \frac{9}{10}$.
}
\end{ex}

\begin{ex}
Rút gọn phân số $\frac{100}{1000}$ về tối giản.
\shortans{1/10}
\loigiai{
$\frac{1}{10}$.
}
\end{ex}

\begin{ex}
Phân số $\frac{13}{17}$ có rút gọn được nữa không? (Được ghi 1, không ghi 0).
\shortans{0}
\loigiai{
Không rút gọn được (0).
}
\end{ex}
